\chapter[Sermon 77. He Shows That Rome Shall Assuredly Fall: And Adds the Causes of Her Fall.]{He Shows That Rome Shall Assuredly Fall: And Adds the Causes of Her Fall.}
\label{sermon:077}
\markright{Sermon 77. He Shows That Rome Shall Assuredly Fall: And Adds the ...}

And after that, I saw an angel come down from Heaven having great power, and the Earth was lightened with his brightness. And he cried mightily with a strong voice, saying, “She has fallen, she has fallen, even great Babylon, and has become the habitation of devils, and the hold of all unclean spirits, and a cage of unclean and hateful birds. For all nations have drunk of the wine of the wrath of her whoredom. And the kings of the earth have committed fornication with her, and her merchants have become rich from the abundance of her pleasures.”

\index{Rome}He pours throughout the 18th chapter the destruction of old and new Rome, also of paganism and Antichristianism, and that with a marvelous abundance and clarity of speech, even so that one would think that you saw all things presently. And he uses also a most godly order. For first, the angel declares the destruction of Rome with most fitting words. Secondly, counsel is given to the godly, how to behave themselves in so great dangers. Then is added the manner of the desolation, that just as Rome has greedily and cruelly spoiled and destroyed other nations, even so it shall befall her also. After this, a lamentation is made, wherein the princes and merchants mourn for the ruin of Rome, where they also recount the riches and pleasures of Rome. Finally, the Apostles and Prophets rejoice at the most just judgment of God. Again, the angel of the Lord cast a millstone into the bottom of the sea, that so the most certain, unrecoverable, and most weighty destruction of Rome might be signified. Whereunto again are annexed the causes of so great evils, and the same finished with the praise and congratulation of all the heavenly dwellers.

\index{Ezekiel (Book of)}And it is fortunate that he imitates the holy prophets of God, of whom two in a manner follow the same pattern, describing the destruction of old Babylon. Isaiah in chapters 13 and 21, and Jeremiah in chapters 50 and 51. And Ezekiel depicts the overthrow of Tyre in chapters 26, 27, and 28. For just as the fate and end of all the ungodly are alike, so the canonical Scripture, in portraying their destruction, agrees well with itself. The Apostles moreover, although they spoke and wrote to the Gentiles in Greek, altered nothing of their natural way of speaking, and even compelled foreign languages to serve the holy, and not the Hebrew to serve pagan languages. For when speaking Greek, they observed the natural phrasing of the Hebrew speech, considering it first, divine, and holy. And though they could speak all languages, they never spoke or wrote any foreign language so that the Hebrew phrasing could not be perceived. Let some therefore beware at this day, that they are not overly fastidious, and follow the purity of the Latin speech so that, in expressing it, they do not fall away from the simplicity of the holy tongue and lose some mysteries. Those who are not stubborn would rather conform themselves to the holy language and learn its phrases than force it against the ear to foreign languages and compel it to serve our delicate ears. Moreover, we have already admonished often what the end and purpose of this treatise is, concerning God’s judgments or punishments. For the truth and justice of God are confirmed, the afflicted receive comfort, and the wicked, and all God’s enemies are made afraid, and so forth.

But when John published these things and prophesied of the destruction of Babylon, all men at that time clearly understood it to signify Rome, because of the recent destruction of Jerusalem and the grievous captivity of the Jews, which had recently occurred under Vespasian. For just as Babylon in times past had vexed the holy city and nation, so now Vespasian, the Roman, did. The godly truly believed these things to be true and that they would undoubtedly come to pass. The ungodly, however, mocked them as foolishness. Their elders had done the same. For when the Prophets also prophesied the destruction of Nineveh, Babylon, and mighty monarchies, they seemed to them to be mad. Nevertheless, just as they had said, so it came to pass. Therefore, the faithful believe God’s oracles, however long they may be deferred, which are prophesied to come, no matter how impossible they may appear to the world. For to God who speaks and wills, nothing is hard.

And intending to demonstrate the downfall of Rome, he prepares his audience and gains credibility for the prophecy, while first of all he shows the author of the oracle or prophecy, the very Angel of God. Indeed, he highly commends this Angel to us, so that we should doubt nothing of the truth of those things which he speaks. For he says how he came from Heaven. From this we gather that those things he brings are divine and celestial, and are said to have great power, lest we should think impossible those things which he says will come to pass. For if God’s Angel is a minister of such power, what may we think the Lord to be, who sent the Angel? One Angel before the walls of Jerusalem killed one hundred eighty-five thousand men of war. One Angel in a night slew all the firstborn of Egypt. Therefore, when the most mighty Angel prophesies the destruction of old and new Rome, we need not doubt that it will utterly perish. Moreover, the Earth was lightened with the glory, that is to say, with the brightness or light of this Angel. For this prophecy is neither obscure nor will it be hidden, but chiefly and most clearly preached throughout the world.

\index{Antichrist}\index{Papacy}Wherefore the same angel cries with all his force. For it behooves these oracles of God, wherein is treated of the glory of God and the salvation of souls, to be preached with loud voices, however the world prohibits and persecutes them. And let those who think that men may be restrained by proclamations, fire, and sword observe that they shall not, with the clearest voice, preach against Antichrist. These fools are deceived. They have fought and contended in this matter for six hundred years and more, and neither could any man, though he raged ever so fiercely, bring this preaching to silence. It breaks out repeatedly and pierces far even at this day throughout the whole world; therefore the glory of this angel is yet, and ever shall be shining and bright, and his voice and preaching most strong, though the Popes [? guts burst].

\index{Jerome, Saint}Now follows the prophecy of Thaungel, some whereof is: Rome shall perish, and no trace of her will remain. He utters this prophetically, as he did also in the 14th Chapter. She is fallen, she is fallen, great Babylon. He said, “She is fallen,” for she shall fall: putting the time past, for the certainty of the thing, for the time to come, whereunto the doubling also appertains. Likewise spake the Prophets. Macrobius marvels at the wonderful brevity of Virgil. And among other things in the first Chapter of the 5th book of Saturnales, will you hear Virgil, says he, speaking with so much brevity, that brevity itself can be no more straightly hampered and drawn together? And fields where Troy was behold how in very few words he has consumed a mighty great city: and has left no ruin at all. Hitherto Macrobius. These things we shall more truly and more rightly apply unto our Prophets, most eloquent in their tongue, and chiefly to Jerome. For what could be thought more brief than that which he said, “She is fallen, she is fallen, great Babylon?” For Jerome both expressed the greatness and majesty of the city, and swallowed it up whole, no ruin at all left, for he signified that both old and new Rome, although it seem stout, invincible, and eternal, yet shall it fall: and so fall, that nothing thereof shall be left. Which shortly after he sets before our eyes more expressly by a certain sign, whilst the Angel taking up a millstone, and casting it into the bottom of the Sea, adds: thus or with such a violence shall Babylon that great city be overthrown, and shall be found no more. Therefore was there never any thing, is, or shall be in the world so mighty or impregnable, which the invincible power of God cannot bring to naught, when he will, and when the fatal hour is come. Old Rome is lost, and that mighty monarchy decayed: there is fall also the superstition and idolatry of the heathens, that has reigned many years: new Rome shall perish also with her Imaginary Empire: the kingdom also of the Pope or Antichrist which has long sotted and plagued the world shall fall, and fade with smoke.

Moreover, by a figurative speech taken from the prophets, he shows the manner of the destruction by consequences: and is become the habitation of devils, and so forth. For so he signifies that it shall be destroyed, that the place which was before much frequented by men, shall now be the habitation of wild beasts and devils, delighting in wilderness, as our Lord also testifies in Matthew 12. And he alluded to the words of the Prophets. Isaiah in chapter 13: Babylon, the beauty of realms, shall be overthrown, as the Lord subverted Sodom and Gomorrah; it shall not be inhabited, but beasts shall there take their rest, and satyrs or harts shall there leap. The same things are repeated also in Jeremiah 50. And in chapter 51, he says: Babylon shall be in heaps, and an habitation for dragons, a water and a hissing, that no man may dwell there. Not unlike things are read in Ezekiel 26 concerning the subversion of Tyre. And that old Rome was destroyed, I showed before: and for the space of forty days and more, inhabited by no man. And that we see it inhabited again, it does not invalidate Christ’s prophecy. For France, Petrarch, an Italian and among the best learned Italians, in a certain epistle to a friend, expounding these words of the Apostle John, among other things, says, “You are verily become such already, for how much better is a wicked man, and of desperate doings, than a devil? Verily you are become the habitation, or rather kingdom of devils: which by their crafts, albeit in human shape, reign in you, and so forth.” Petrarch lived and wrote these things about two hundred years ago. And in another certain epistle, speaking of old and new Babylon, he says, “She was, of all others, worst, and at that time most filthy; and this now is no city, but an house of fiends and sprites, and to be short, the sink of all sin and shame, and that hell of the living, signified long before by the mouth of David, than it was founded or known.” And the same again: whatever you have read of Babylon in Assyria or Egypt, whatever you have read of the four Labyrinths or Mazes, finally whatever you have read of the way to hell, of the dark words there and lakes of fire and brimstone, compared to this hell, it is a fable: here is that proud and terrible Nimrod; here is Semiramis with her quiver; here is unmerciful Minos; here is Rhadamanthus; here is Cerberus devouring all things; here is Pasiphae put to the bull, among a cruel kind, as Virgil says, a young of double shape, Minotaurus by name, a monstrous progeny of unlawful lust. Finally, here you may see whatever is confused, whatever is black, whatever is or may be feigned horrible and ugly, and so forth. These things he has written, and many other more like these, in other epistles. But what do you think he would write now, if he saw the court of Rome at this day? Which is doubtless in many ways more corrupt than it was then. Briefly, John signifies after the sentence of Christ, king and judge, that Rome, both old and new, together with paganism and Antichristianism, shall perish utterly, and never be restored again.

The causes I have rehearsed once or twice, he repeats and emphasizes: ungodliness, idolatry, and the deception of all people and nations, whom they have compelled by torments to receive impiety. Where cruelty, tyranny, and bloodshed also hold sway. I spoke of the wine of whoredom before in chapter 14 and elsewhere, so it needs not to be repeated again with tediousness. And hereunto is added another new cause, and the merchants of the Earth who grew rich from the power or influence of her pleasures. And he said, of the power of pleasure: for through immeasurable, mighty, and insatiable lust they grew rich. For as Rome abounded with spoils, which it had greedily taken from all nations and brought to Rome, they were given to all kinds of riot and wantonness. Therefore the masters of luxury, the dividers of delicate pleasures, and the merchants of the most precious wares, repairing thither, always found that they would be entertained and honored, and were made rich by the luxurious and riotous life of the Romans. Therefore the Apostle notes an incredible pursuit of most sumptuous riot, in food, drink, apparel, building, in pampering and cherishing of the body. The Roman church also of our time, struck with the same rage both in Italy and elsewhere, accumulates excessive riches through riotous living. This is seen chiefly in those spiritual fathers, bishops, and abbots, and in the entire Roman clergy. But God never allowed riot and tyranny to go unpunished in any nation. Therefore Babylon has fallen, and therefore the church of Rome will fall too. Therefore let private men also love temperance, and abstain from riot and pride. To the Lord be glory.
